\chapter{\texorpdfstring{\(L^1\)}{L1} Fourier 变换}
\label{chap:l1-fourier-transform}

周期函数的频率是整数，定义在整条直线上的函数则需要连续频率。把周期长度
送向无穷，离散系数的求和形式转为频率变量上的积分，得到 Fourier 变换。
本章在 \(\R^d\) 上从 \(L^1\) 理论开始，因为绝对可积性允许直接使用 Fubini
定理和控制收敛定理。采用约定
\begin{equation}\label{eq:fourier-transform-convention}
 \widehat f(\xi)=\int_{\R^d}f(x)e^{-2\pi i x\cdot\xi}\,dx.
\end{equation}
这一约定使卷积公式和反演公式都不带额外常数；微分公式中出现 \(2\pi i\xi\)。

\section{定义与基本估计}

若 \(f\in L^1(\R^d)\)，则积分\eqref{eq:fourier-transform-convention}绝对收敛，
且
\begin{equation}\label{eq:fourier-lone-bound}
 |\widehat f(\xi)|\le\|f\|_1.
\end{equation}
因此 Fourier 变换是从 \(L^1(\R^d)\) 到有界函数空间的有界线性映射，算子
范数为一。范数恰为一可用非负积分为一的函数在 \(\xi=0\) 处取等号说明。

\begin{theorem}\label{thm:fourier-lone-continuous-czero}
若 \(f\in L^1(\R^d)\)，则 \(\widehat f\) 一致连续且在无穷远处趋于零。
换言之，Fourier 变换把 \(L^1(\R^d)\) 映入 \(C_0(\R^d)\)。
\end{theorem}
\begin{proof}
对 \(h\to0\)，
\[
 |\widehat f(\xi+h)-\widehat f(\xi)|
 \le\int|f(x)|\,|e^{-2\pi ix\cdot h}-1|\,dx.
\]
右侧与 \(\xi\) 无关；被积函数逐点趋零并由 \(2|f|\) 控制，故一致连续。

再证无穷远消失。平移在 \(L^1\) 中连续。给定 \(\xi\ne0\)，取
\(h=\xi/(2|\xi|^2)\)，则 \(e^{-2\pi ih\cdot\xi}=-1\)。换元给
\[
 2\widehat f(\xi)
 =\int_{\R^d}\bigl(f(x)-f(x-h)\bigr)e^{-2\pi ix\cdot\xi}\,dx,
\]
所以 \(2|\widehat f(\xi)|\le\|f-\tau_hf\|_1\)。当
\(|\xi|\to\infty\) 时 \(h\to0\)，右侧趋零。
\end{proof}

这里没有假定 \(f\) 非负。若另有 \(f\ge0\)，则
\(|\widehat f(\xi)|\le\widehat f(0)\)，并且变换满足章末习题中的正定性条件。
这些额外限制来自非负性，不能移用于任意实值或复值的 \(L^1\) 函数。

\section{平移、调制与伸缩}

令 \((\tau_af)(x)=f(x-a)\)、\((M_bf)(x)=e^{2\pi ib\cdot x}f(x)\)。
直接换元得到
\begin{align}
 \widehat{\tau_af}(\xi)&=e^{-2\pi ia\cdot\xi}\widehat f(\xi),\label{eq:fourier-translation}\\
 \widehat{M_bf}(\xi)&=\widehat f(\xi-b).\label{eq:fourier-modulation}
\end{align}
空间中的平移对应频率中的相位，空间中的振荡因子对应频率平移。

若 \(A\) 是可逆实矩阵，令 \(g(x)=f(Ax)\)。变量代换给
\begin{equation}\label{eq:fourier-linear-change}
 \widehat g(\xi)=|\det A|^{-1}\widehat f(A^{-T}\xi).
\end{equation}
特别地，对 \(\varepsilon>0\)，\(f_\varepsilon(x)=\varepsilon^{-d}f(x/\varepsilon)\) 满足
\(\widehat{f_\varepsilon}(\xi)=\widehat f(\varepsilon\xi)\)。

\section{乘法与微分}

多重指标记作 \(\alpha=(\alpha_1,\ldots,\alpha_d)\in\N^d\)，
\(|\alpha|=\sum\alpha_j\)、\(x^\alpha=\prod x_j^{\alpha_j}\)。
写 \(\beta\le\alpha\) 表示每个坐标上都有 \(\beta_j\le\alpha_j\)。
若对所有 \(\beta\le\alpha\) 都有 \(x^\beta f\in L^1\)，可在积分号下逐次求导：
\begin{equation}\label{eq:fourier-moment-derivative}
 \partial_\xi^\alpha\widehat f(\xi)
 =(-2\pi i)^{|\alpha|}\widehat{x^\alpha f}(\xi).
\end{equation}
事实上，一次差商中的因子满足
\[
 \left|\frac{e^{-2\pi ihx_j}-1}{h}\right|\le2\pi|x_j|.
\]
因而 \(x_jf\in L^1\) 允许使用控制收敛，得到一阶公式。再把同一论证用于
\(x^\beta f\)，便逐次得到所需偏导；每一步的连续性由
\(L^1\) Fourier 变换的连续性保证。一个便于使用的充分条件是
\((1+|x|^{|\alpha|})f\in L^1\)。

多元情形不能省略中间的矩条件。例如在 \(\R^2\) 上取
\(f=\mathbf1_{\{x>1,\ 0<y<x^{-2}\}}\)，则
\[
 \int f=1,\qquad \int xyf=\frac12\int_1^\infty x^{-3}\,dx<\infty,
 \qquad \int xf=\int_1^\infty x^{-1}\,dx=\infty.
\]
仅知道最高的混合矩可积，并不提供逐次求导所需的各个控制函数。

反过来，若 \(f\in C_c^1(\R^d)\)，分部积分的边界项为零，从而
\begin{equation}\label{eq:fourier-derivative-multiplier}
 \widehat{\partial_jf}(\xi)=2\pi i\xi_j\widehat f(\xi).
\end{equation}
\chapterref{chap:distributions-weak-derivatives}建立开集上的弱导数后，
这一公式还将推广到 \(f,\partial_jf\in L^1\) 的情形。届时先用光滑截断和
磨光在 \(L^1\) 中同时逼近函数与该方向的弱导数，再取极限。
本章后续证明只使用已经得到的紧支撑 \(C^1\) 版本。
微分提高频率权重，乘以坐标则在频率中求导；这对偶关系是后面用 Fourier
变换处理微分方程的起点。

\section{卷积代数}

对 \(f,g\in L^1(\R^d)\)，定义
\[
 (f*g)(x)=\int_{\R^d}f(x-y)g(y)\,dy.
\]
Tonelli 定理给
\[
 \|f*g\|_1\le\|f\|_1\|g\|_1,
\]
故卷积几乎处处有定义且仍在 \(L^1\)。换元与 Fubini 定理给
\begin{theorem}[卷积定理]\label{thm:fourier-convolution}
\[
 \widehat{f*g}=\widehat f\,\widehat g.
\]
\end{theorem}
\begin{proof}
二重绝对积分为
\(\iint|f(x-y)g(y)|\,dy\,dx=\|f\|_1\|g\|_1\)。故换序合法，
再作 \(z=x-y\)，得
\[
 \widehat{f*g}(\xi)
 =\int g(y)e^{-2\pi iy\cdot\xi}\,dy
   \int f(z)e^{-2\pi iz\cdot\xi}\,dz.
\]
\end{proof}
乘积公式是本章反演定理成立后的推论，此处先说明其条件。若 \(f,g\in L^1\) 且
\(\widehat f,\widehat g\in L^1\)，反演给出二者的有界连续代表，故
\(fg\in L^1\)。将 \(g\) 的反演积分代入，并以
\(\|f\|_1\|\widehat g\|_1\) 控制二重绝对积分，得到
\[
 \widehat{fg}(\xi)=\int\widehat f(\xi-\eta)\widehat g(\eta)\,d\eta.
\]
这里的乘积可积来自反演的附加假设；仅凭 \(f,g\in L^1\) 不能断言它成立。

卷积交换、结合，并对平移不变。没有 \(L^1\) 函数充当卷积单位：若
\(e*f=f\) 对所有 \(f\) 成立，对每个固定 \(\xi\) 取
\(f(x)=e^{2\pi ix\cdot\xi}h(x)\)，其中 \(h\in L^1\) 且 \(\int h=1\)，
则 \(\widehat f(\xi)=1\)。卷积定理给 \(\widehat e(\xi)=1\)，这与
Riemann--Lebesgue 引理矛盾。Dirac 质量将在分布空间中补上这一单位。

\section{Gaussian 函数与逼近恒等}

令
\[
 \gamma(x)=e^{-\pi|x|^2},\qquad
 \gamma_\varepsilon(x)=\varepsilon^{-d}\gamma(x/\varepsilon).
\]
先核实归一化常数。令 \(I=\int_\R e^{-\pi x^2}dx\)，尾部与指数衰减比较
可知 \(0<I<\infty\)。Tonelli 定理、圆盘上的极坐标变换和单调收敛给
\[
 I^2=\int_{\R^2}e^{-\pi(x^2+y^2)}\,dx\,dy
 =\lim_{R\to\infty}2\pi\int_0^R r e^{-\pi r^2}\,dr
 =\lim_{R\to\infty}(1-e^{-\pi R^2})=1.
\]
因为 \(I>0\)，得到 \(I=1\)。再按坐标使用 Tonelli，
\(\int_{\R^d}\gamma=I^d=1\)。这里及下文均取 \(\varepsilon>0\)。

\begin{proposition}\label{prop:gaussian-fourier}
\(\widehat\gamma=\gamma\)，从而
\(\widehat{\gamma_\varepsilon}(\xi)=e^{-\pi\varepsilon^2|\xi|^2}\)。
\end{proposition}
\begin{proof}
先在一维。积分号下求导并分部积分，记 \(G=\widehat\gamma\)，得
\[
 G'(\xi)=-2\pi i\int x e^{-\pi x^2}e^{-2\pi ix\xi}dx
 =-2\pi\xi G(\xi).
\]
因此 \((e^{\pi\xi^2}G(\xi))'=0\)，结合 \(G(0)=1\) 得
\(G(\xi)=e^{-\pi\xi^2}\)。
高维函数与指数核均按坐标分解，Fubini 定理给乘积结论；伸缩公式给第二式。
\end{proof}

族 \((\gamma_\varepsilon)\) 是逼近恒等：对 \(1\le p<\infty\)，
\[
 \|f*\gamma_\varepsilon-f\|_p^p
 \le\int\gamma(y)\|\tau_{\varepsilon y}f-f\|_p^p\,dy\longrightarrow0.
\]
不等式来自概率测度 \(\gamma(y)\,dy\) 下的 Jensen 不等式，再对 \(x\)
使用 Tonelli 定理。对 \(p=1\)，它就是积分三角不等式。
同样的估计先用于 \(f*\gamma_\varepsilon\)，保证卷积几乎处处存在。
对固定 \(y\) 使用平移连续性，并以
\(2^p\|f\|_p^p\gamma(y)\) 控制，即得极限。
这里证明的是范数收敛，尚未从中断言整个族逐点收敛。

\section{Fourier 反演}

\begin{theorem}[Fourier 反演定理]\label{thm:fourier-inversion-lone}
设 \(f\in L^1(\R^d)\) 且 \(\widehat f\in L^1(\R^d)\)。则在 \(f\) 的每个
Lebesgue 点 \(x\)，
\[
 f(x)=\int_{\R^d}\widehat f(\xi)e^{2\pi ix\cdot\xi}\,d\xi.
\]
因而该右侧给出 \(f\) 的一个连续代表；等式至少几乎处处成立。
\end{theorem}
\begin{proof}
对 \(\varepsilon>0\)，绝对可积性允许使用 Fubini：
\begin{align*}
 \int \widehat f(\xi)e^{2\pi ix\cdot\xi}
       e^{-\pi\varepsilon^2|\xi|^2}\,d\xi
 &=\int f(y)\left[\int
 e^{2\pi i(x-y)\cdot\xi}e^{-\pi\varepsilon^2|\xi|^2}\,d\xi\right]dy\\
 &=(f*\gamma_\varepsilon)(x).
\end{align*}
记反演积分为 \(F(x)\)。由本章的连续性定理，\(F\in C_0(\R^d)\)。
左侧与 \(F\) 的差一致地不超过
\[
 \int|\widehat f(\xi)|\,|e^{-\pi\varepsilon^2|\xi|^2}-1|\,d\xi
 \longrightarrow0.
\]
另一方面，右侧在 \(L^1\) 中趋于 \(f\)。从 \(\varepsilon=1/n\) 中抽取
一个使右侧几乎处处收敛的子列，就得到 \(F=f\) 几乎处处。

最后，称 \(x\) 为这里的 Lebesgue 点，是指
\[
 \frac1{|B(x,r)|}\int_{B(x,r)}|f(y)-f(x)|\,dy\longrightarrow0.
\]
在这样的点，利用 \(F=f\) 几乎处处，有
\[
 |F(x)-f(x)|\le
 \sup_{y\in B(x,r)}|F(x)-F(y)|
 +\frac1{|B(x,r)|}\int_{B(x,r)}|f(y)-f(x)|\,dy.
\]
让 \(r\downarrow0\)，两项都趋零。这一步只使用连续性和 Lebesgue 点的定义，
不需要预先调用高维 Lebesgue 微分定理。
\end{proof}

反演假设中的 \(\widehat f\in L^1\) 不能由 \(f\in L^1\) 自动推出。例如
区间指标函数的变换按 \(1/|\xi|\) 衰减，在一维并不可积。对这类函数仍可用
对称截断、Cesàro 平均或 \(L^2\) 理论讨论反演，但不能套用本定理。

\begin{corollary}[唯一性]\label{cor:fourier-lone-uniqueness}
若 \(f,g\in L^1(\R^d)\) 且 \(\widehat f=\widehat g\)，则 \(f=g\) 几乎处处。
\end{corollary}
\begin{proof}
令 \(h=f-g\)。有 \(h\in L^1\)，而 \(\widehat h=0\) 本身也属于 \(L^1\)。
反演定理直接给 \(h=0\) 几乎处处。
\end{proof}

\section*{习题}
\addcontentsline{toc}{section}{习题}
\subsection*{A组}
\begin{enumerate}[label=\textbf{\thechapter.\arabic*.}]
\exerciseitem{fourier-transform-indicator-interval}设 \(a>0\)，计算 \(\mathbf1_{[-a,a]}\) 的 Fourier 变换，并判断其变换是否属于 \(L^1(\R)\)。
\exerciseitem{fourier-transform-translation-modulation}逐项证明平移与调制公式\eqref{eq:fourier-translation}—\eqref{eq:fourier-modulation}。
\exerciseitem{fourier-transform-dilation}证明可逆线性变量代换公式\eqref{eq:fourier-linear-change}。
\exerciseitem{fourier-transform-continuity}用控制收敛证明 \(\widehat f\) 一致连续，并指出估计为何与基点 \(\xi\) 无关。
\end{enumerate}
\subsection*{B组}
\begin{enumerate}[label=\textbf{\thechapter.\arabic*.},resume]
\exerciseitem{fourier-transform-moments}在对所有 \(\beta\le\alpha\) 都有 \(x^\beta f\in L^1\) 的条件下证明公式\eqref{eq:fourier-moment-derivative}。
\exerciseitem{fourier-transform-derivative}对 \(f\in C_c^1(\R^d)\) 证明公式\eqref{eq:fourier-derivative-multiplier}；学完弱导数后，再用稠密逼近推广到 \(f,\partial_jf\in L^1\) 的情形。
\exerciseitem{fourier-convolution-associative}证明 \(L^1(\R^d)\) 中卷积交换、结合，且满足 Young 不等式。
\exerciseitem{fourier-no-lone-identity}补全“\(L^1\) 卷积代数无单位”的论证。
\exerciseitem{fourier-gaussian-approximate-identity}证明 Gaussian 族在 \(L^p\) 中构成逼近恒等，\(1\le p<\infty\)。
\end{enumerate}
\subsection*{C组}
\begin{enumerate}[label=\textbf{\thechapter.\arabic*.},resume]
\exerciseitem{fourier-inversion-continuous-representative}证明若 \(f,\widehat f\in L^1\)，则 \(f\) 几乎处处等于一个属于 \(C_0\) 的函数。
\exerciseitem{fourier-heat-kernel}设 \(t>0\)，求 \(e^{-4\pi^2t|\xi|^2}\) 的逆 Fourier 变换，并验证它关于 \(t\) 满足热方程。
\optionalexerciseitem{fourier-positive-definite}若 \(f\ge0\) 且 \(f\in L^1\)，证明 \(\widehat f\) 是正定函数；再说明正定性比“连续且趋零”多给出了什么限制。这里称 \(F\) 正定，是指对任意有限个 \(\xi_j\in\R^d\)、\(c_j\in\C\)，都有 \(\sum_{j,k}c_j\overline{c_k}F(\xi_j-\xi_k)\ge0\)。
\end{enumerate}
